\HLASourcePage{177}{176}
\thispagestyle{hla-body}
\begin{hlafixedpage}
  \HLAFixedBlock{90.000}{69.786}{432.000}{\rmfamily\raggedright}{9.96}{12.00}{%
    Logical time values shall be used to represent a federate’s time position and timestamp values. TRADT%
  }
  \HLAFixedBlock{90.000}{81.785}{432.000}{\rmfamily\raggedright}{9.96}{12.00}{%
    specific methods shall be used to compare the relative order of one logical time with another and to compare%
  }
  \HLAFixedBlock{90.000}{93.785}{432.000}{\rmfamily\raggedright}{9.96}{12.00}{%
    a logical time value with the initial value and with the final value. Similar methods shall be used to add and%
  }
  \HLAFixedBlock{90.000}{105.785}{432.000}{\rmfamily\raggedright}{9.96}{12.00}{%
    subtract a value to and from of a TRADT instance of logical time, which enforce compliance with the%
  }
  \HLAFixedBlock{90.000}{117.785}{432.000}{\rmfamily\raggedright}{9.96}{12.00}{%
    required logical time properties.%
  }
  \HLAFixedBlock{90.000}{141.785}{432.000}{\rmfamily\raggedright}{9.96}{12.00}{%
    Logical time interval (LTI) values shall be used to represent lookahead values and intervals of logical time.%
  }
  \HLAFixedBlock{90.000}{153.784}{432.000}{\rmfamily\raggedright}{9.96}{12.00}{%
    There is no concept of a “final” or “largest possible” LTI value, but LTIs shall not be negative. However,%
  }
  \HLAFixedBlock{90.001}{165.784}{431.999}{\rmfamily\raggedright}{9.96}{12.00}{%
    TRADT specific methods shall be used to test if an LTI value is zero in addition to those which add or%
  }
  \HLAFixedBlock{90.001}{177.784}{431.999}{\rmfamily\raggedright}{9.96}{12.00}{%
    subtract a value to and from an LTI value.%
  }
  \HLAFixedBlock{90.001}{201.784}{431.999}{\rmfamily\raggedright}{9.96}{12.00}{%
    The arithmetic operations allowed on logical time values shall be restricted. These restrictions permit%
  }
  \HLAFixedBlock{90.002}{213.783}{431.998}{\rmfamily\raggedright}{9.96}{12.00}{%
    TRADTs to be generalized, expressing properties not always present in integer or floating-point numbers.%
  }
  \HLAFixedBlock{90.002}{225.783}{431.998}{\rmfamily\raggedright}{9.96}{12.00}{%
    Addition and subtraction shall only be defined when the addend or subtrahend is an LTI. While an LTI can%
  }
  \HLAFixedBlock{90.002}{237.783}{431.998}{\rmfamily\raggedright}{9.96}{12.00}{%
    be added to an LTI to produce a new LTI, logical time values can only have an LTI added to them or%
  }
  \HLAFixedBlock{90.002}{249.783}{431.998}{\rmfamily\raggedright}{9.96}{12.00}{%
    subtracted from them. Comparison operations: less than, less than or equal to, equal to, equal to or greater%
  }
  \HLAFixedBlock{90.002}{261.783}{431.998}{\rmfamily\raggedright}{9.96}{12.00}{%
    than, greater than, and not equal to shall be defined upon logical time values and LTIs, but not between%
  }
  \HLAFixedBlock{90.002}{273.783}{431.998}{\rmfamily\raggedright}{9.96}{12.00}{%
    them. Logical time values may include positive and negative values between the initial and final values. No%
  }
  \HLAFixedBlock{90.002}{285.782}{431.998}{\rmfamily\raggedright}{9.96}{12.00}{%
    other operations shall be permitted.%
  }
  \HLAFixedBlock{90.002}{309.782}{431.998}{\rmfamily\raggedright}{9.96}{12.00}{%
    Logical time has the following notable values:%
  }
  \HLAFixedBlock{100.022}{324.782}{421.978}{\rmfamily\raggedright}{9.96}{12.00}{%
    a)%
  }
  \HLAFixedBlock{121.972}{324.782}{400.028}{\rmfamily\raggedright}{9.96}{12.00}{%
    initial time value--The point on the time axis for which no other point occurs before it.%
  }
  \HLAFixedBlock{100.022}{339.781}{421.978}{\rmfamily\raggedright}{9.96}{12.00}{%
    b)%
  }
  \HLAFixedBlock{121.982}{339.781}{400.018}{\rmfamily\raggedright}{9.96}{12.00}{%
    final time value--The point on the time axis for which no other points occurs after it.%
  }
  \HLAFixedBlock{90.002}{363.781}{431.998}{\rmfamily\raggedright}{9.96}{12.00}{%
    Similarly, LTI has the following notable values:%
  }
  \HLAFixedBlock{100.022}{378.781}{421.978}{\rmfamily\raggedright}{9.96}{12.00}{%
    a)%
  }
  \HLAFixedBlock{121.965}{378.781}{400.035}{\rmfamily\raggedright}{9.96}{12.00}{%
    Zero time interval--The distance between two logical time values that are equal. When this logical%
  }
  \HLAFixedBlock{121.982}{390.781}{400.018}{\rmfamily\raggedright}{9.96}{12.00}{%
    time interval is added to a logical time T1 the result is a logical time T2 == T1.%
  }
  \HLAFixedBlock{100.022}{405.780}{421.978}{\rmfamily\raggedright}{9.96}{12.00}{%
    b)%
  }
  \HLAFixedBlock{121.982}{405.780}{400.018}{\rmfamily\raggedright}{9.96}{12.00}{%
    Epsilon time interval--The distance between any two consecutive logical times T1 and T2 (i.e. the%
  }
  \HLAFixedBlock{121.982}{417.780}{400.018}{\rmfamily\raggedright}{9.96}{12.00}{%
    set of logical times between T1 and T2 is empty). For some representations of local time (e.g.%
  }
  \HLAFixedBlock{121.982}{429.780}{400.018}{\rmfamily\raggedright}{9.96}{12.00}{%
    floating point), the distance may not be the same for all values of T1 and T2.%
  }
  \HLAFixedBlock{90.002}{453.780}{431.998}{\rmfamily\raggedright}{9.96}{12.00}{%
    LTIs have unusual properties when implemented upon floating-point numbers. As floating-point numbers%
  }
  \HLAFixedBlock{90.002}{465.779}{431.998}{\rmfamily\raggedright}{9.96}{12.00}{%
    become larger the distance between adjacent discrete values becomes larger. Thus, the epsilon time interval%
  }
  \HLAFixedBlock{90.001}{477.779}{431.999}{\rmfamily\raggedright}{9.96}{12.00}{%
    depends on the logical time to which it is being added or subtracted. When the addition or subtraction%
  }
  \HLAFixedBlock{90.001}{489.779}{431.999}{\rmfamily\raggedright}{9.96}{12.00}{%
    operations are provided a non-zero LTI which is smaller than the epsilon time interval for the logical time%
  }
  \HLAFixedBlock{90.001}{501.779}{431.999}{\rmfamily\raggedright}{9.96}{12.00}{%
    value to which they are added or subtracted, the result shall be adding or subtracting the epsilon time interval%
  }
  \HLAFixedBlock{90.001}{513.779}{431.999}{\rmfamily\raggedright}{9.96}{12.00}{%
    instead.%
  }
  \HLAFixedBlock{90.001}{537.778}{431.999}{\rmfamily\raggedright}{9.96}{12.00}{%
    Logical time values have unusual properties when they approach their initial and final values. Operations%
  }
  \HLAFixedBlock{90.001}{549.778}{431.999}{\rmfamily\raggedright}{9.96}{12.00}{%
    that would otherwise result in a value larger than the final value shall result in the final value instead.%
  }
  \HLAFixedBlock{90.001}{561.778}{431.999}{\rmfamily\raggedright}{9.96}{12.00}{%
    Similarly, operations that would otherwise result in a value lower than the initial value shall result in the%
  }
  \HLAFixedBlock{90.001}{573.778}{431.999}{\rmfamily\raggedright}{9.96}{12.00}{%
    initial value instead. Operations that would result in an LTI less than zero result in a zero interval instead.%
  }
  \HLAFixedBlock{90.001}{597.777}{431.999}{\rmfamily\raggedright}{9.96}{12.00}{%
    When a logical time value is equal to the final value it has a “sticky” property with respect to subtraction.%
  }
  \HLAFixedBlock{90.000}{609.777}{432.000}{\rmfamily\raggedright}{9.96}{12.00}{%
    The addition or subtraction of a logical time interval to or from a logical time having the final value shall%
  }
  \HLAFixedBlock{90.000}{621.777}{432.000}{\rmfamily\raggedright}{9.96}{12.00}{%
    result in the final value.%
  }
  \HLAFixedBlock{90.000}{645.776}{432.000}{\rmfamily\raggedright}{9.96}{12.00}{%
    Given initial time Ti, final time Tf, logical time T between Ti and Tf, nonzero intervals I \& I1 \& I2, Zero%
  }
  \HLAFixedBlock{90.000}{657.776}{432.000}{\rmfamily\raggedright}{9.96}{12.00}{%
    time interval Iz, and epsilon time interval Ie, all implementations of logical time shall confirm that the%
  }
  \HLAFixedBlock{90.000}{669.776}{432.000}{\rmfamily\raggedright}{9.96}{12.00}{%
    following statements are true.%
  }
\end{hlafixedpage}
\clearpage
